\documentclass[10pt,a4paper]{amsart}
\usepackage[margin=19mm]{geometry}
\usepackage{amsmath,amssymb,mathtools}
\usepackage[scr=rsfs,cal=euler]{mathalfa}
\usepackage{xfrac}
\usepackage[unicode,hidelinks,bookmarks=false]{hyperref}
\newcommand{\ie}{i.e.\ }
\newcommand{\cf}{cf.\ }
\newcommand{\resp}{respectively\ }
\newcommand{\Subsection}{\subsection}
\providecommand{\nobreakdash}{\nobreak\hbox{-}}
\setlength{\emergencystretch}{2em}
\multlinegap0pt
\usepackage{bm}
\usepackage{bbm}
\usepackage{dsfont}
\usepackage{xspace}
\usepackage{cancel}
\usepackage{mathtools}
\usepackage{amsthm}
\newtheorem{theorem}{Theorem}
\newtheorem{proposition}[theorem]{Proposition}
\newcommand{\VH}{\;\!\!V\!(\!H\;\!\!)\!}
\newcommand{\lcw}{\operatorname{lcw}}
\title[Linear clique-width and modular decomposition]{Linear clique-width and modular decomposition}
\author{Robert Brignall, Michal Opler,, Vincent Vatter}
\date{Source: arXiv:2602.22089v1 (CC BY 4.0)}
\begin{document}
\maketitle

\noindent\emph{Source and licence.} Linear clique-width and modular decomposition. Version arXiv:2602.22089v1 (CC BY 4.0), distributed under the Creative Commons Attribution 4.0 International licence, as recorded in the arXiv licence metadata for this exact version and shown on the abstract page https://arxiv.org/abs/2602.22089v1. Full rights notice: licenses/arxiv-2602.22089-CC-BY-4.0.txt.\par\smallskip

\noindent\textit{Editorial scope.} Counted unit: the Proposition whose source label is \texttt{prop:main} in the article (source0.tex lines 425--432, source label \texttt{prop:main}), delivered together with the article's own complete original proof of it (source0.tex lines 434--501, one \texttt{proof} environment of 68 lines). No mathematics is written, repaired or re-derived: statement and proof are the article's own text line for line. Required context retained uncounted: prop:lcw:inflation:comp (lines 365--371); prop:lcw:two:modules (lines 390--404). Every definition, display and label of those lines is retained unchanged. Omitted from this package and not redistributed: 1--364, 372--389, 405--424, 433--433, 502--561. Nothing inside a retained excerpt is omitted or altered: every retained body is delivered exactly as the article's lines are written, so no omission receipt of any kind accompanies this package. Citations: the retained excerpts carry no citation command at all, so no bibliography entry of the article is transcribed and no reference is redistributed. Reference adaptation: none was needed and none was made; every reference inside the delivered unit resolves inside it, so no unresolved reference or citation is produced. Printed slips kept literal: no printed word, symbol, hypothesis or number of the counted statement or of its proof was altered. Layout and class: the article's own class is \texttt{article} with packages including fontenc, amsmath, amsthm, amssymb, latexsym, enumitem, calc, xspace, tikz, titlesec, caption, natbib, footmisc, color; the wrapper is the lane's pinned \texttt{amsart} editorial wrapper at 10pt on A4, which restates the theorem-like environments the retained text needs and adds the article's own macro definitions that the retained text needs (\texttt{\textbackslash VH}, \texttt{\textbackslash lcw}), with extra wrapper packages (bm, bbm, dsfont, xspace, cancel, mathtools). The delivered numbering is the unit's own. Assets: no retained span contains a figure environment, a graphics-inclusion command, a PDF-page inclusion, a table or a TikZ picture; no image file is copied into this package, the package declares no asset dependency and no image file is redistributed with it. Licence: version arXiv:2602.22089v1 of this article is distributed under the Creative Commons Attribution 4.0 International licence, as recorded in the arXiv licence metadata for this exact version and shown on the abstract page https://arxiv.org/abs/2602.22089v1. A full rights notice is delivered at licenses/arxiv-2602.22089-CC-BY-4.0.txt. Characters: every retained excerpt is pure ASCII, so the delivered engine, XeLaTeX with its default Latin Modern text font, reports no missing character.
\begin{proposition}
\label{prop:lcw:inflation:comp}
If $G=H[\{G_v : v\in V(H)\}]$ where $H$ is complete or anti-complete, then
\[
	\lcw(G)\le 1+\max_{v\in \VH} \lcw(G_v).
\]
\end{proposition}

\begin{proposition}
\label{prop:lcw:two:modules}
Suppose that $G=H[G_v : v\in V(H)]$ where $H$ has at least two vertices. Then either there is a vertex $x$ of $H$ for which
\[
	\lcw(G_x)
	=
	\lcw(G),
\]
or there are distinct vertices $x$ and $y$ of $H$ for which
\[
	\lcw(G_x),\ \lcw(G_y)
	\ge
	\lcw(G)-\lcw(H)-1.
\]
\end{proposition}

\begin{proposition}
\label{prop:main}
Let $G$ be a graph for which the linear clique-width of every prime induced subgraph is at most $m$, and let $t,s\ge 1$. If
\[
	\lcw(G)\ge (m+2)(t+s),
\]
then $G$ contains $Q_t$ or $\overline{Q}_s$ as an induced subgraph.
\end{proposition}

\begin{proof}
We proceed by induction on the number of vertices of $G$. We may assume that $m\ge 2$, as otherwise $G$ cannot contain $K_2$, so $G$ is anti-complete and cannot satisfy the bound in the proposition for any $t,s\ge 1$. In particular, $G$ has at least two vertices, establishing the base case of our induction.

We may also assume that $t,s\ge 3$. Cases with $t=1$ or $s=1$ are trivial, since $Q_1=\overline{Q}_1=K_1$ is an induced subgraph of every nonempty graph. For $t=2$, note that a graph not containing $Q_2=K_2\uplus K_1$ as an induced subgraph is complete multipartite, with linear clique-width at most~$2$. Since $m\ge 2$ and $s\ge 1$, we have $(m+2)(2+s)\ge 12>2$, so every graph satisfying the bound contains $Q_2\le Q_t$. Dually, a graph not containing $\overline{Q}_2=P_3$ is a disjoint union of cliques, with linear clique-width at most~$3$. Since $m\ge 2$ and $t\ge 1$, we have $(m+2)(t+2)\ge 12>3$, so every graph satisfying our bound also contains $\overline{Q}_2\le\overline{Q}_s$. We make this reduction because $t,s\ge 3$ ensures that~$Q_{t-1}$ is disconnected and $\overline{Q}_{s-1}$ is connected, properties we use in the argument below.

%
%
%

Now suppose that $G$ satisfies the hypotheses of the proposition with $m\ge 2$ and $t,s\ge 3$ and that the result holds for all graphs with fewer vertices. Apply the first stage of the modular decomposition to write
\[
	G=H[G_v : v\in V(H)]
\]
where $H$ is complete, anti-complete, or prime on at least four vertices, and in the first two cases the modules $G_v$ are the co-components or connected components of $G$, respectively. Let $x$ be a vertex of $H$ for which $\lcw(G_x)$ is maximal amongst all modules $G_v$. If $\lcw(G_x)=\lcw(G)$, then we are done by induction because $G_x$ itself satisfies the hypotheses and therefore contains either $Q_t$ or $\overline{Q}_s$.

Otherwise, Proposition~\ref{prop:lcw:two:modules} provides a vertex $y\neq x$ of $H$ such that
\[
	\lcw(G_x),\ \lcw(G_y)
	\ge
	\lcw(G)-\lcw(H)-1.
\]
Since $\lcw(H)\le m$ (as $H$ is an induced subgraph of $G$ that is either complete, anti-complete, or prime), we have
\[
	\lcw(G_x),\ \lcw(G_y)
	\ge
	\lcw(G)-m-1
	\ge
	(m+2)(t+s-1)+1.
\]
Since $G_x$ and $G_y$ have fewer vertices than $G$, the inductive hypothesis applies to them for all values of $t$ and $s$. Applying it with $(t-1,s)$ and $(t,s-1)$, we conclude that each of $G_x$ and $G_y$ contains~$Q_{t-1}$ or $\overline{Q}_s$, and each also contains $Q_t$ or $\overline{Q}_{s-1}$. If either contains $Q_t$ or $\overline{Q}_s$, then so does $G$ and we are done. We may therefore assume that both $G_x$ and~$G_y$ contain $Q_{t-1}$ and $\overline{Q}_{s-1}$.

%
%
%

First suppose that $H$ is prime on four or more vertices. Because $\{x,y\}$ does not form a module in~$H$, there is some vertex $z\in V(H)$ adjacent to precisely one of $x$ or $y$. If $x$ and $y$ are adjacent in~$H$, then~$G$ contains $\overline{Q}_s=(K_1\uplus \overline{Q}_{s-1})\ast\overline{Q}_{s-1}$, while if $x$ and $y$ are nonadjacent in $H$, then~$G$ contains $Q_t=(K_1\ast Q_{t-1})\uplus Q_{t-1}$, and in either case we are done.

%
%
%

It remains to consider the case where $H$ is complete or anti-complete. We treat the anti-complete case first, and then explain the modifications needed for the complete case.

Suppose that $H$ is anti-complete, so $G$ is the disjoint union of its connected components, two of which are $G_x$ and $G_y$. By Proposition~\ref{prop:lcw:inflation:comp}, $\lcw(G_x)\ge\lcw(G)-1\ge (m+2)(t+s)-1$. Since $G_x$ is connected, its skeleton is either complete or prime on at least four vertices.

If $G_x$ is the inflation of a complete graph (a join of co-components), then since $Q_{t-1}$ is disconnected, it lies entirely within a single co-component of $G_x$, and any vertex from another co-component is adjacent to all of $Q_{t-1}$. Thus $G_x$ contains $K_1\ast Q_{t-1}$, which together with the copy of $Q_{t-1}$ in $G_y$ gives $Q_t=(K_1\ast Q_{t-1})\uplus Q_{t-1}$ in $G$.

%
%

If $G_x$ is the inflation of a prime graph $H'$ on four or more vertices, we apply Proposition~\ref{prop:lcw:two:modules} to $G_x$. Let $a$ be a vertex of $H'$ maximizing $\lcw(G'_a)$. If $\lcw(G'_a)=\lcw(G_x)$, then $G'_a$ has fewer vertices than $G$ and satisfies the hypotheses of the proposition, so we are done by induction. Otherwise, Proposition~\ref{prop:lcw:two:modules} provides a vertex $b\neq a$ of $H'$ with
\[
	\lcw(G'_a),\ \lcw(G'_b)
	\ge
	\lcw(G_x)-m-1
	\ge
	(m+2)(t+s-1).
\]
Applying the inductive hypothesis to $G'_a$ and $G'_b$ as we did with $G_x$ and $G_y$ earlier, we may assume that they each contain $Q_{t-1}$ and $\overline{Q}_{s-1}$.

Since $H'$ is prime, $\{a,b\}$ is not a module, so there exists a vertex $c\in V(H')$ adjacent to precisely one of $a$ or $b$. Say $c$ is adjacent to $a$. Any vertex of $G'_c$ is then adjacent to all of $Q_{t-1}$ inside $G'_a$, giving a copy of $K_1\ast Q_{t-1}$ inside $G_x$. Combined with the copy of $Q_{t-1}$ in $G_y$ (a separate component of $G$), we obtain $Q_t=(K_1\ast Q_{t-1})\uplus Q_{t-1}$ in $G$, completing the proof in this case.

%
%
%

Now suppose that $H$ is complete, so $G$ is the join of its co-components, two of which are $G_x$ and $G_y$. By Proposition~\ref{prop:lcw:inflation:comp}, $\lcw(G_x)\ge (m+2)(t+s)-1$ as before. Since $G_x$ is co-connected, its skeleton is either anti-complete or prime on at least four vertices. If $G_x$ is the inflation of an anti-complete graph (a disjoint union of connected components), then $\overline{Q}_{s-1}$ is connected (since $s\ge 3$), so it lies within a single component of $G_x$, and any vertex from another component is nonadjacent to all of~$\overline{Q}_{s-1}$. Thus $G_x$ contains $K_1\uplus\overline{Q}_{s-1}$, which together with the copy of~$\overline{Q}_{s-1}$ in~$G_y$ (joined to~$G_x$ since~$H$ is complete) gives $\overline{Q}_s=(K_1\uplus\overline{Q}_{s-1})\ast\overline{Q}_{s-1}$ in~$G$. If instead $G_x$ has a prime skeleton, we find two large modules inside~$G_x$ each containing~$Q_{t-1}$ and~$\overline{Q}_{s-1}$, exactly as in the anti-complete case. The prime skeleton of $G_x$ again provides a vertex $c$ adjacent to precisely one of them, say $c$ is adjacent to $a$ but not to $b$. A vertex of $G'_c$ is then nonadjacent to all of $\overline{Q}_{s-1}$ inside $G'_b$, giving $K_1\uplus\overline{Q}_{s-1}$ inside $G_x$. Together with the copy of $\overline{Q}_{s-1}$ in $G_y$ (joined to $G_x$ since $H$ is complete), we obtain $\overline{Q}_s=(K_1\uplus\overline{Q}_{s-1})\ast\overline{Q}_{s-1}$ in~$G$. This completes the proof.
\end{proof}

\end{document}
